%% !TEX root = ../manual_gaolv5.tex
% The abstract and the keywords (unnumbered).
%%
%% Copyright (c) 2026 ENSTA, France
%%
%% By: Jordan NININ (created 2026-11-06)
%%
%% GAOL is a software distributed WITHOUT ANY WARRANTY. Read the associated
%% COPYING.LIB file for information.
\chapter*{Abstract}
\thispagestyle{plain}

GAOL is a C++ library for arithmetic on floating-point intervals. Its
intervals enclose real numbers between two doubles, and every operation
returns an interval that contains the exact result, whatever the rounding
errors of the computer: this is the containment property, on which interval
algorithms rest to prove what they compute. GAOL also provides the
\emph{relational} versions of its operations, the reverse functions of IEEE
1788-2015, which interval constraint programming needs. GAOL v5, which this
manual describes, continues GAOL as Fr\'ed\'eric Goualard developed it, and
is written so that it compiles and computes right on every system it can:
Linux, macOS and Windows, on x86, x86-64, 32-bit and 64-bit ARM, POWER
(ppc64le), s390x and RISC-V (riscv64).

\begin{itemize}
\item[$\bullet$] Every elementary function that CORE-MATH provides is bounded
  with it, which is correctly rounded in the rounding direction in effect: the
  value computed at a bound, rounding upward, is the upper bound, and the
  double below it the lower one, the tightest bounds. The integer powers, the
  roots of an order other than 2 and 3 and \code{compoundm1}, which GAOL computes itself, are the tightest
  bounds or within a few doubles of them (Chapter~\ref{chap:accuracy}). The
  bounds are the same on every architecture and with every compiler, and no
  mathematical library is to be installed or linked: CORE-MATH is compiled
  into GAOL.
\item[$\bullet$] GAOL v5 follows IEEE 1788-2015 in the operations GAOL
  provides: the empty set, the infinite bounds, the domains of the functions,
  the comparisons and the interval literals. The namespace
  \code{gaol\_ieee1788} gives each operation of the standard the name and the
  argument order the standard gives it, and the tightness of every operation
  is documented, as the standard requires.
\item[$\bullet$] The rounding direction is the business of GAOL: each
  operation that computes a bound sets the direction upward when it is not,
  whichever direction the calling code left, and GAOL can be built to restore
  the direction of the program after each operation; the operations of the
  AVX-512 path, which that build takes where the processor has the
  instructions, set no direction, their instructions carrying it.
\item[$\bullet$] Numbers are read exactly, and intervals written in decimal
  are rounded outward, whatever the C library; written in hexadecimal, an
  interval is read back as the same doubles.
\item[$\bullet$] Three builds, CMake, autotools and meson, configure GAOL the
  same way; the tests compare the bounds GAOL computes with their exact
  values, and a continuous integration builds and tests GAOL on Linux, macOS
  and Windows, with the compilers and on the processors it supports.
\end{itemize}

\medskip
\noindent
This manual is both a user's guide and a reference. It presents the notions
of interval arithmetic GAOL relies on, the installation and the use of the
library, each of its operations with examples whose outputs were produced by
GAOL v5, the names of the standard, the accuracy of the operations, the way
the intervals are computed and measured, the tests and the continuous
integration, and what a program written for GAOL~4 has to change.

\bigskip
\noindent\textbf{Keywords:} interval arithmetic, floating-point arithmetic,
IEEE 1788-2015, correctly rounded elementary functions, CORE-MATH, reverse
functions, constraint programming, C++.
